\documentclass[a4paper]{article}
% Español (México) con XeLaTeX: fontspec para fuentes Unicode, polyglossia
% para la división silábica y los nombres del documento en la variante mexicana.
\usepackage{fontspec}
\usepackage{polyglossia}
\setdefaultlanguage[variant=mexican]{spanish}
% Los nombres chinos van como «pinyin (original)»: xeCJK da a los caracteres
% Han su propia fuente; sin ella XeLaTeX los omite con un aviso.
% FandolSong no cubre algunos caracteres de nombres (U+73FA); WenQuanYi Zen Hei sí.
\usepackage[AutoFallBack=true]{xeCJK}
\setCJKmainfont{FandolSong-Regular.otf}
\setCJKfallbackfamilyfont{\CJKrmdefault}{WenQuanYi Zen Hei}
% polyglossia no traduce los nombres de listings.
\AtBeginDocument{\providecommand{\lstlistingname}{}\renewcommand{\lstlistingname}{Listado}%
  \providecommand{\lstlistlistingname}{}\renewcommand{\lstlistlistingname}{Índice de listados}}
\usepackage{amsmath, amssymb}
\usepackage{graphicx}
\usepackage[margin=2.5cm]{geometry}
\usepackage[most]{tcolorbox}
\usepackage{etoolbox}
\usepackage{listings}
\usepackage{hyperref}
\usepackage{booktabs}
\usepackage{subcaption}
\usepackage{float}

\newtcolorbox{knowledgebox}[1]{enhanced, colback=blue!5!white, colframe=blue!75!black, colbacktitle=blue!75!black, coltitle=white, fonttitle=\bfseries, title={#1}, attach boxed title to top left={yshift=-2mm, xshift=2mm}, boxrule=1pt, sharp corners}
\newtcolorbox{importantbox}[1]{enhanced, colback=yellow!10!white, colframe=yellow!80!black, colbacktitle=yellow!80!black, coltitle=black, fonttitle=\bfseries, title={#1}, sharp corners}
\newtcolorbox{warningbox}[1]{enhanced, colback=red!5!white, colframe=red!75!black, colbacktitle=red!75!black, coltitle=white, fonttitle=\bfseries, title={#1}, sharp corners}

\lstset{language=Python, basicstyle=\ttfamily\small, keywordstyle=\color{blue}, stringstyle=\color{red!60!black}, commentstyle=\color{green!60!black}, breaklines=true, frame=single, numbers=left, numberstyle=\tiny\color{gray}, captionpos=b, extendedchars=true}

\newcommand{\notetitle}{{[}LLM Agents F24{]} Enterprise Trends for Generative AI and Building Successful Agents — Burak Gokturk}
\newcommand{\noteauthors}{Elaboradas a partir de materiales públicos del curso}
\newcommand{\notedate}{30 de septiembre de 2024}
\newcommand{\videochannel}{Berkeley RDI}
\newcommand{\videopublishdate}{30 de septiembre de 2024}
\newcommand{\videoduration}{aproximadamente 70 minutos}
\newcommand{\videourl}{https://www.youtube.com/live/Sy1psHS3w3I}
\newcommand{\videocoverpath}{cover.jpg}
\newcommand{\repourl}{https://github.com/hqhq1025/ai-course-notes}


% Footer with repo info
\usepackage{fancyhdr}
\pagestyle{fancy}
\fancyhf{}
\fancyfoot[C]{\thepage}
\fancyfoot[R]{\footnotesize\color{gray}\href{\repourl}{AI Course Notes}}
\renewcommand{\headrulewidth}{0pt}
\renewcommand{\footrulewidth}{0pt}

\begin{document}
\begin{titlepage}\centering
{\Large Notas del curso\par}\vspace{1.2cm}
{\huge\bfseries \notetitle\par}\vspace{0.8cm}
{\large \noteauthors\par}\vspace{0.3cm}
{\large \notedate\par}\vspace{1.2cm}
\ifdefempty{\videocoverpath}{{\small Favor de completar la ruta de la imagen de portada.\par}}{\includegraphics[width=0.82\textwidth,height=0.45\textheight,keepaspectratio]{\videocoverpath}\par}
\vfill
\begin{tcolorbox}[width=0.9\textwidth, colback=black!2!white, colframe=black!60, sharp corners]
\textbf{Autor o canal del video}: \videochannel\par
\textbf{Fecha de publicación}: \videopublishdate\par
\textbf{Duración del video}: \videoduration\par
\ifdefempty{\videourl}{}{\textbf{Enlace del video}: \href{\videourl}{\nolinkurl{\videourl}}\par}\par
\textbf{Página del proyecto}: \href{\repourl}{\nolinkurl{\repourl}}\par
\end{tcolorbox}
\end{titlepage}

\tableofcontents\newpage

\section{Introducción: tendencias de la GenAI desde la perspectiva empresarial}

Burak Gokturk es vicepresidente de Ingeniería de Google. Esta clase no profundiza en detalles algorítmicos, sino que examina las tendencias de desarrollo de la GenAI desde la perspectiva macro del \textbf{mercado empresarial}, para ofrecer a estudiantes y emprendedores una orientación práctica.

\begin{knowledgebox}{El cambio de perspectiva del doctorado a la industria}
Burak Gokturk abre con su propia experiencia de doctorado: en aquella época se consideraba que las redes neuronales «no eran tan buenas como las SVM», y entrenar los datos de tomografía computarizada de 50 pacientes en 42 segundos ya se percibía como demasiado lento. Hoy en día, entrenar los modelos grandes toma semanas enteras. Una preocupación interesante es que los ciclos de entrenamiento tan prolongados han hecho más lenta la velocidad con la que los investigadores adquieren intuición.
\end{knowledgebox}

\subsection{Resumen de la sección}

La velocidad de desarrollo de la IA supera con creces lo esperado, y tanto la academia como la industria están viviendo un cambio de paradigma.

\section{Tendencias clave de la GenAI empresarial}

\subsection{Democratización de la IA}

\begin{importantbox}{Los desarrolladores sustituyen a los expertos en IA}
Antes, construir aplicaciones de IA requería científicos de datos y expertos en ML; ahora cualquier desarrollador —incluso un estudiante de secundaria— puede construir aplicaciones de IA usando la API de un LLM. Esto amplió el grupo de desarrolladores de IA en más de 100 veces, y es la causa fundamental del crecimiento explosivo de las aplicaciones de IA.
\end{importantbox}

Detrás de la democratización de la IA hay dos fuerzas: la primera es que cada vez más modelos base abren sus API; la segunda es que las herramientas de bajo código o sin código encapsulan el diseño de prompts, la limpieza de datos y los procesos de Ops en flujos de trabajo visuales. Juntas, estas dos fuerzas trasladaron las aplicaciones de IA de un puñado de data scientist a un gran número de ingenieros de software y equipos de negocio.

\subsection{Disminución de la demanda de datos}

\begin{itemize}
    \item Antes, las empresas necesitaban meses o incluso años para recopilar datos etiquetados
    \item Los modelos grandes ya contienen un enorme conocimiento de preentrenamiento, por lo que la cantidad de datos necesaria para el fine-tuning se redujo drásticamente
    \item Las empresas pueden poner en marcha un proyecto de IA en cuestión de días, no de años
\end{itemize}

Reducir la demanda de datos no significa que se pueda omitir la gobernanza de datos. Las empresas aún necesitan establecer estrategias de «catálogo de datos», «puntuación de calidad de datos» y «acceso diferenciado» para evitar que el modelo absorba contenido no conforme durante el entrenamiento. Muchas empresas optan por validar primero con datos públicos y después hacer el ajuste fino con un lote pequeño de datos de alta calidad.

\subsection{Iteración rápida de modelos}

\begin{itemize}
    \item Cada pocas semanas se publica un nuevo modelo, y las clasificaciones cambian constantemente
    \item Las empresas se enfocan en la \textbf{plataforma} y no en un modelo individual: necesitan una infraestructura que permita cambiar de modelo con flexibilidad
    \item Hace 18 meses solo había un modelo de punta (GPT-4); ahora hay varias opciones equiparables
\end{itemize}

La iteración frecuente también trae consigo problemas de gestión de versiones. Una de las características de Vertex AI es que permite definir varias versiones (intercambiables) dentro del mismo pipeline de Agent y registrar los resultados de experimentos A/B. Así, incluso si alguna versión presenta un regress, es posible revertirla rápidamente.

\subsection{Convergencia multimodal}

\begin{itemize}
    \item Los modelos unimodales están siendo sustituidos por modelos multimodales (Gemini admite de forma nativa texto, imagen, audio y video)
    \item La demanda empresarial de entradas multimodales (gráficos en documentos, análisis de contenido de video) crece día con día
\end{itemize}

La convergencia multimodal implica atender simultáneamente problemas de sincronización de datos, alineación y resolución. Gemini construye un embedding space unificado entre voz, imagen y texto, de modo que un mismo flujo de entrada puede ser analizado a la vez por varios tokenizers, lo que evita rupturas entre modalidades.

\subsection{Descenso continuo de costo y latencia}

\begin{itemize}
    \item El costo de las llamadas a la API ha bajado alrededor de 1.5 órdenes de magnitud
    \item La reducción de la latencia hace posibles más aplicaciones en tiempo real
    \item Los modelos MoE dispersos (que solo activan una parte de los parámetros) reducen aún más el costo de la inferencia
    \item Hallazgo interesante: ejecutar ciertos modelos de API resulta más económico que ejecutar modelos de código abierto equivalentes (porque se ahorran los costos de operación)
\end{itemize}

Tras la baja de costos, las empresas prefieren la ruta de «algoritmo + hardware heterogéneo»: activar GPU de baja latencia en las horas pico, habilitar TPU/CPU más económicas en las horas de menor demanda, e incluso hacer offload de una parte de las solicitudes a un datacenter regional. Este enfoque exige que la capa de orchestration pueda percibir las características reales de cada solicitud.

\subsection{Resumen de la sección}

Tendencias clave de la GenAI empresarial: democratización de la IA, disminución del umbral de datos, iteración rápida de modelos, convergencia multimodal y descenso continuo de costos.

\section{Línea de tiempo de la evolución técnica de GenAI}

\subsection{2018-2021: la era del foundation model}

El mayor avance de esta etapa es el preentrenamiento + fine-tuning de Transformer. Burak menciona el GTM de esa época: primero se usan BERT y GPT-2 para resolver tareas de NLU/NLG, y después se abre el embedding a la Search empresarial y al document understanding. La práctica típica de esta etapa consiste en building an "AI co-pilot" for narrow tasks, y se enfatiza que «la preparación de datos, el feature engineering y el knowledge graph» son la verdadera ventaja competitiva.

\subsection{2022-2023: Agent y multi-turn}
El concepto de Agent comienza a popularizarse en 2022: el LLM ya no solo produce texto, sino que orchestrating tools. Tras la aparición de frameworks como Vertex AI Agent Builder y LangChain, los ingenieros pueden armar en pocos días un flujo plan→tool→feedback. El problema central de esta etapa es «hard-coded workflow + manual monitoring», por lo que Burak insiste en hacer metrics-first design.

En esta etapa es habitual que el repository incluya: prompt templates, tool manifest, evaluation harness. Burak menciona en particular que hay que documentar el sample input/output, estimated latency y failure modes de cada tool en un doc para que el planner los consulte. De lo contrario, el agent puede provocar fácilmente un cascade failure al interpretar mal un tool call.

\subsection{Después de 2024: tiempo real, control y cumplimiento}

A partir de 2024, el interés de las empresas se desplaza hacia el tiempo real y el cumplimiento normativo. El multi-modal de Gemini, el agent scheduler de Vertex y el RSP de Anthropic exigen todos un nivel más alto de «model governance». Las palabras clave de la nueva etapa son «trazabilidad», el «policy-as-code» visualizado y un «failure analysis» riguroso.

Una tendencia clave es que «legal hold + audit trail» se está convirtiendo gradualmente en baseline. El equipo de cumplimiento exige que cada decision path pueda reproducirse, e incluso que el user prompt y el tool call se encadenen en una timeline para su revisión. Muchas empresas conectan esto a un data lake y lo vinculan con el pipeline de observabilidad.

\subsection{Resumen de la sección}

Al repasar la línea de tiempo se observa que, del preentrenamiento al Agent y de ahí al cumplimiento normativo, el reto de la IA empresarial se desplaza continuamente de la capability hacia el control, lo que implica que los equipos deben iterar sin cesar sus procesos y su estructura organizacional.

\section{Prácticas típicas de GenAI empresarial}

\subsection{Asistentes inteligentes en los servicios financieros}

La ruta habitual de GenAI en la banca digital consiste en integrar el LLM en el servicio al cliente, la gestión patrimonial y los procesos internos de cumplimiento normativo. El patrón típico es el siguiente:

\begin{itemize}
    \item \textbf{Entrada}: el usuario pregunta por asesoría de inversión / reporta una transacción sospechosa
    \item \textbf{Búsqueda}: recopila en tiempo real datos de mercado y normativa regulatoria
    \item \textbf{LLM}: genera una respuesta explicativa más una recomendación de acción
    \item \textbf{Herramientas}: llena automáticamente los reportes y genera borradores de contratos
\end{itemize}

\begin{knowledgebox}{El ciclo cerrado de tres capas de la IA financiera}
1. Frontend: el LLM genera una respuesta explicativa<br>
2. Capa intermedia: la búsqueda/recuperación dinámica garantiza que la información esté actualizada<br>
3. Backend: herramientas de automatización de bajo código ejecutan las instrucciones<br>
Entre las tres capas debe garantizarse la trazabilidad (¿por qué se recomienda esta estrategia?) y la prevención de riesgos (rechazar operaciones sensibles).
\end{knowledgebox}

\subsection{El GenAI Agent en la manufactura y la cadena de suministro}

Las empresas manufactureras usan el GenAI Agent para conectar el MES, el ERP y las herramientas de visualización de la cadena de suministro:

\begin{itemize}
    \item El ingeniero sube el CAD junto con la descripción del caso de uso
    \item El Agent lee el BOM, los registros de prueba y el historial de fallas, y genera un plan de diagnóstico de fallas
    \item Activa directamente una orden de trabajo o notifica la planeación de compras
\end{itemize}

\begin{warningbox}{El escenario de prueba debe reproducir la realidad}
Muchos agentes se desempeñan bien en el entorno de prueba, pero fallan en cuanto entran a un taller real debido a la diferencia en la distribución de los datos (bitácoras frente a sensores en tiempo real). Es necesario recopilar feedback durante la operación real y reproducirlo para el modelo.
\end{warningbox}

\subsection{Resumen de la sección}

Las prácticas de los distintos sectores siguen el mismo patrón de ciclo cerrado: la IA aporta comprensión y generación, la búsqueda y los datos garantizan la exactitud factual, y las herramientas y procesos completan la ejecución; toda la cadena requiere trazabilidad y control.

\section{Búsqueda + LLM: la aplicación estrella de las empresas}

\begin{importantbox}{Grounding with Search}
Casi todos los casos de uso empresariales combinan un LLM con la búsqueda:
\begin{itemize}
    \item El conocimiento del LLM tiene un límite de vigencia temporal: se necesita la búsqueda para obtener información actualizada
    \item El LLM produce alucinaciones: se necesita la búsqueda para verificar los hechos
    \item El LLM es bueno para razonar pero carece de referencias: la búsqueda aporta citas precisas de las fuentes
\end{itemize}
Esto hace que RAG (generación aumentada por recuperación) se haya vuelto la arquitectura estándar de las aplicaciones empresariales de IA generativa.
\end{importantbox}

Desde el punto de vista de la ingeniería, el pipeline de RAG debe resolver tres problemas clave: primero, la latencia que trae consigo el crecimiento del índice vectorial; segundo, las anomalías de recuperación (que se devuelvan documentos obsoletos o erróneos); tercero, la estrategia de integración entre el LLM y los resultados recuperados. Burak comentó que algunos equipos agregan una capa de «fact filter» antes de los resultados de la recuperación: primero se le entregan los resultados a un smaller model para hacer un reordenamiento fino, y después se le entregan al LLM grande para la generación.

\subsection{Resumen de la sección}

Búsqueda + LLM es el patrón de aplicación de IA generativa más común en las empresas, y resuelve los tres puntos críticos de la vigencia temporal, las alucinaciones y las referencias.

\section{Módulos clave para construir un AI Agent}

\subsection{Arquitectura de la plataforma Vertex AI}

El diseño de Vertex AI de Google refleja las necesidades comunes de una plataforma de IA empresarial:

\begin{itemize}
    \item \textbf{Capa de chips}: flexibilidad para elegir GPU/TPU (actualmente hay una escasez global severa de chips)
    \item \textbf{Capa de modelos}: una interfaz unificada para modelos propios, de terceros y de código abierto
    \item \textbf{Capa de herramientas}: Model Builder y Agent Builder ofrecen capacidades de personalización y orquestación
    \item \textbf{Capa de aplicaciones}: aplicaciones verticales como búsqueda, conversación y análisis de datos
\end{itemize}

\subsection{Diferenciación de Gemini}

\begin{itemize}
    \item Multimodalidad nativa (mezcla varias modalidades desde el inicio del entrenamiento)
    \item Ventana de contexto extremadamente grande (en los experimentos ya alcanza 10 millones de tokens)
    \item Prueba de Needle-in-a-haystack: recuperación precisa de información en contextos extremadamente largos
\end{itemize}

\subsection{Módulos de Planner y Memory}

Dentro de un Agent normalmente hay tres componentes que colaboran: planner, executor y memory. El flujo típico es:

\begin{itemize}
    \item \textbf{Planner}: descompone la intención del usuario en pasos (por ejemplo, búsqueda, llamada a herramientas, resumen)
    \item \textbf{Memory}: registra cada conversación, consulta y salida de herramientas, y ofrece context re-use
    \item \textbf{Executor}: se encarga de invocar herramientas específicas, API o entornos de ejecución personalizados
\end{itemize}

\begin{importantbox}{El reto de sincronización entre planner y memory}
Una vez que el planner define los pasos, memory necesita devolver rápidamente los resultados de ejecuciones anteriores. Si memory no es lo bastante detallada (solo conserva la ronda más reciente), el planner volverá a llamar a las herramientas de forma repetida; si memory es demasiado grande, esto provocará latency. La solución es introducir compressed memory + recall hints, es decir, mantener en memory un query summary como referencia para el planner.
\end{importantbox}
\begin{importantbox}{División de funciones entre Planner y Memory}
Memory le proporciona al planner una caché de contexto; el planner, factores de decisión (el fracaso de la ronda anterior, la tasa de éxito de las herramientas); executor se ejecuta bajo un guardrail de seguridad. Esta arquitectura de tres capas le permite al agent mantener coherencia en tareas de contexto largo que abarcan varias herramientas.
\end{importantbox}

\subsection{Tool Orchestration y Workflow}

Para que el Agent coordine varias herramientas (búsqueda, bases de datos, API), se necesita un flujo de trabajo y una tool specification claros:

\begin{itemize}
    \item Definir tool intent, inputs y outputs, para evitar model hallucinating tool usage
    \item Usar `tool descriptor` para validar los parámetros de llamada del LLM (tipo, rango)
    \item Insertar `pre-conditions` y `post-conditions` en el workflow
\end{itemize}

\begin{knowledgebox}{Diseño del workflow como state machine}
Un agent complejo a menudo modela el workflow como una state machine: cada estado representa una fase de la tarea (comprensión, planning, execution, confirmation); la condición que dispara la transition proviene del juicio del planner. Esto permite simular distintas ramas con unit test durante la fase de desarrollo, y mejora la system reliability.
\end{knowledgebox}
\begin{knowledgebox}{El valor del tool descriptor}
Cuando el modelo necesita invocar `fetch\_legislation()` o `run\_cad\_simulation()`, el tool descriptor ofrece schema, authority y expected latency, para evitar que el agent llame a ciegas o que se agote el tiempo de espera.
\end{knowledgebox}

\subsection{Resumen de la sección}

Una plataforma de IA empresarial necesita ofrecer flexibilidad de elección en cada una de las capas: chips, modelos, herramientas y aplicaciones.

\section{Observabilidad y control de seguridad del AI Agent}

\subsection{Indicadores de monitoreo y estrategia de alertas}

Una plataforma de Agent madura primero debe definir indicadores observables y después generar alertas automáticas en cuanto aparezcan «señales»:

\begin{itemize}
    \item \textbf{Consistencia semántica}: Historical Reference vs. current response divergence
    \item \textbf{Tasa de alucinación del modelo}: revisar aleatoriamente las salidas para ver si se introdujeron hechos falsos
    \item \textbf{Errores en la llamada a herramientas}: fallas en la ejecución de comandos / tiempos de espera agotados en llamadas a la API
    \item \textbf{Explosión del contexto}: una expansión repentina de la ventana de contexto indica un memory leak o una prompt injection
    \item \textbf{Número de correcciones del usuario}: interrupciones frecuentes del usuario indican que el agent malinterpretó el intent
\end{itemize}

\begin{knowledgebox}{Cuanto antes se capturen las signals, más fácil es corregir el rumbo}
Muchos equipos solo se enteran de que el agent tiene un problema cuando el usuario se queja. Las señales anómalas tempranas suelen esconderse en un latency spike, en tiempos de espera agotados en la llamada a herramientas o en «contenido repetido». Establecer un esquema transparente de «metric + sampling reforzado + human-in-loop review» permite escalar la respuesta antes de que el agent falle.
\end{knowledgebox}

Burak subraya que el pipeline de datos de instrumentation debe mantener un costo operativo bajo: para la consistencia semántica o la hallucination, solo se muestrea entre el 1 y el 2\% de las solicitudes, y la sensibilidad se preserva mediante la revisión del squad. Recolectar en exceso genera costos y también puede provocar fatiga de alertas.

\subsection{Gobernanza de seguridad y diseño de guardrail}

El comportamiento del Agent no puede depender únicamente del prompt del modelo; también depende de guardrails a nivel de sistema:

\begin{itemize}
    \item \textbf{Policy Layer}: intercepta operaciones sensibles antes de ejecutarlas (eliminar archivos, programar tareas)
    \item \textbf{Tool Firewall}: limita el acceso del modelo a las herramientas a una lista blanca
    \item \textbf{Human-in-the-loop level}: decide dinámicamente, según el riesgo de la tarea, si se necesita aprobación humana
    \item \textbf{Audit Trail}: registra cada decisión, cada llamada a herramientas y cada confirmación del usuario
\end{itemize}

\begin{warningbox}{Sin audit trail no hay manera de exigir responsabilidades}
En cuanto un agent provoca un error o un incidente de seguridad, la falta de un registro de operaciones significa que no se puede reconstruir la cadena de decisiones. Esto hace que el sistema de gobernanza deje de funcionar e incluso afecta la auditoría de cumplimiento. Tanto Anthropic como Vertex AI subrayan que cada tool call debe ser trazable y reversible.
\end{warningbox}

\subsection{Resumen de la sección}

Solo mediante indicadores cuantitativos, cadenas de alerta y cadenas de responsabilidad puede una empresa reaccionar con rapidez cuando el agent presenta anomalías, y mantener así los límites de seguridad al mismo tiempo que crece su capacidad.

\section{Despliegue y operación de AI Agent}

\subsection{Trade-off entre costo y latencia}

Desplegar un Agent no es solo la selección del modelo, sino también resource orchestration:

\begin{itemize}
    \item \textbf{Arquitectura de inferencia}: chat + tool loop vs. batch prompt, distintas estrategias tienen un impacto enorme en el latency
    \item \textbf{Escalonamiento de recursos}: el path caliente usa low-latency GPU, el path frío usa cheaper CPU
    \item \textbf{Auto-scaling}: escala automáticamente según el request rate, para evitar overprovisioning
    \item \textbf{Multi-model routing}: las tareas de bajo riesgo usan un smaller model, las de alto riesgo usan Sonnet/Opus
\end{itemize}

\begin{importantbox}{No basta con mirar el rendimiento, hay que mirar la end-to-end experience}
El éxito de un Agent no es simplemente la velocidad de inferencia, sino el tiempo total de finalización de la tarea: respuesta del LLM→llamada a la API→posprocesamiento→confirmación del usuario. Un retraso en cualquier paso puede arruinar la experiencia del usuario. El Agent Builder de Vertex AI, mediante un tablero de pipeline, reúne el latency de cada componente de manera concreta en el dashboard.
\end{importantbox}

La lección aquí es que el tamaño del Pod y el prompt size forman dos dimensiones que se presionan mutuamente. Para una token-heavy application, se puede dividir el pipeline en un «first pass de grano grueso» y un «second pass fino», minimizando el costo de model call sin dejar de garantizar el latency.

\subsection{Ciclo de datos y retroalimentación del usuario}

Mejorar continuamente el Agent requiere reutilizar en ciclo el real user data:

\begin{itemize}
    \item \textbf{Structured feedback}: insertar botones de «¿se resolvió el problema?» y «¿cumple con las normas?»
    \item \textbf{Failure tagging}: etiquetar los casos de fallo (hallucination, tool failure, safety violation)
    \item \textbf{Retraining loop}: llevar los failure de alto valor a fine-tuning/LoRA
    \item \textbf{Prompt atlas}: versionar cada modificación de prompt, para facilitar la reversión
\end{itemize}

\begin{knowledgebox}{El Feedback loop da resultados antes que el «reentrenamiento»}
Muchas empresas primero usan el user feedback para ajustar el prompt, y después deciden si reentrenan. Esta estrategia de «primero ajustar la interacción, después hacer un retrain a gran escala» reduce el costo de iteración y permite entregar un Agent confiable más rápido.
\end{knowledgebox}

En Vertex AI, el feedback loop normalmente corre en un `data labeling project` independiente: después de enviar los casos de fallo al labeler para su anotación, se dispara automáticamente el retrain pipeline. Este ciclo automatizado garantiza que el model loop se mantenga fresh en todo momento.

\subsection{Resumen de la sección}

La operación del Agent se enfoca en el costo de latencia, la proporción de recursos y el ciclo de datos; solo después de estandarizar estos elementos el equipo puede llevar las capacidades de laboratorio al entorno de producción.

\section{Consejos para estudiantes y emprendedores}

\begin{itemize}
    \item Dirección de investigación: ¿cómo personalizar un modelo grande para un caso de uso específico? ¿cómo sustituir la búsqueda tradicional con contexto largo?
    \item Lo que más necesita el campo de la IA es \textbf{creatividad}: no solo capacidad técnica, sino también pensamiento interdisciplinario
    \item Se recomienda estudiar humanidades, arte u otras ciencias además de la IA, para cultivar el pensamiento creativo
\end{itemize}

\subsection{Resumen de la sección}

Lo que más necesitan los estudiantes y emprendedores es convertir la curiosidad técnica en proyectos interdisciplinarios: integrar producto, experiencia y valor social para construir servicios de GenAI que realmente se puedan implementar, en lugar de solo escribir artículos.

\section{Evaluación y laboratorio del agente LLM}

\subsection{Métricas cuantitativas de evaluación}

Antes de desplegar un agente, la empresa necesita construir un conjunto de métricas cuantitativas:

\begin{itemize}
    \item \textbf{Task Success Rate}: tasa de cumplimiento del objetivo (por ejemplo, organizar automáticamente un informe y obtener su aprobación)
    \item \textbf{Force Stop Rate}: proporción de conversaciones que el usuario interrumpe de forma activa
    \item \textbf{Tool Error Rate}: errores en la API o en las herramientas que provocan un workflow failure
    \item \textbf{Response Divergence}: la consistencia de las salidas entre distintas versiones del modelo
\end{itemize}

Estas métricas suelen organizarse por niveles de SLA: las Critical Task usan 99.9\% success, mientras que las Exploratory Task permiten un umbral más bajo.

\subsection{Laboratorios Human-in-the-loop}

Burak menciona que las empresas establecen un «Agent Lab»: testing con un grupo reducido de usuarios frente a escenarios reales. El proceso es el siguiente:

\begin{enumerate}
    \item Se simulan tareas reales y se deja que el agente se ejecute en un sistema tan realista como sea posible
    \item Se asigna un reviewer humano para hacer audit output de manera periódica y etiquetarlo
    \item Los resultados del label se convierten en un dashboard que dispara el retrain o el prompt tune
\end{enumerate}

\begin{knowledgebox}{El valor organizacional del Agent Lab}
El Agent Lab es el bridge between research and product: el equipo de investigación puede ver failure reales, el equipo de producto puede probar la nueva versión, y el equipo de gobernanza puede verificar si los guardrails de seguridad son eficaces.
\end{knowledgebox}

\subsection{Resumen de la sección}

El sistema de evaluación se compone de métricas cuantitativas y de laboratorios humanos, lo que garantiza que el agente sea eficiente y que, en caso de falla, se pueda localizar rápidamente el root cause.

\section{Gobernanza y alineación con la regulación}

\subsection{Requisitos de política y cumplimiento}

Los GenAI/Agent en las empresas suelen tocar campos sensibles como finanzas, salud y gobierno. Es necesario cumplir simultáneamente con los siguientes requisitos regulatorios:

\begin{itemize}
    \item \textbf{Soberanía de datos}: los datos deben procesarse en la región designada y no pueden filtrarse a través de fronteras
    \item \textbf{Transparencia}: el usuario debe poder ver por qué el agent hizo determinada recomendación
    \item \textbf{Auditoría de seguridad}: todas las llamadas a tool deben registrarse y quedar disponibles para revisión del equipo de cumplimiento
    \item \textbf{Detección de sesgos}: el agent no puede producir salidas discriminatorias por raza o género
\end{itemize}

\subsection{Combinación de tecnología y gobernanza}

El cumplimiento no depende solo de documentos, sino de codificar el guardrail como code:

\begin{itemize}
    \item policy-as-code: escribir las reglas como policy ejecutable (por ejemplo, Open Policy Agent)
    \item simulation sandbox: cada vez que se modifica el prompt, primero se ejecuta una prueba en el sandbox, para asegurar que no se viole la policy
    \item kill-switch: cuando se detecta un ataque o un abuso, el agent puede entrar automáticamente en un mode seguro
\end{itemize}

\begin{warningbox}{La regulación no es un checkbox, sino un diálogo continuo}
Los organismos reguladores plantean continuamente nuevos requisitos conforme a las aplicaciones reales. Las empresas necesitan integrar el regulator view en el daily standup: cuando aparece una new capability, primero se evalúa si activa una nueva policy, y solo después se pone en producción.
\end{warningbox}

\subsection{Resumen de la sección}

La gobernanza requiere colaboración entre organizaciones: legal, cumplimiento, seguridad de AI e ingeniería construyen juntos el policy-as-code, el interruptor de apagado y un rastro de auditoría explicable.

\section{Implicaciones organizacionales y de inversión}

\subsection{Estructura de talento y cultura}

Burak enfatiza en particular la importancia de la densidad de talento y de una cultura de iteración: organizar los squad multifuncionales en una tríada de product + safety + research. Dentro de cada squad se necesitan cuatro tipos de roles —prompt engineer, search engineer, ops engineer y policy analyst— para asegurar que el agent pipeline no solo pueda ejecutarse, sino también administrarse.

La experiencia de Anthropic en 2023 es que la densidad de talento importa más que la cantidad. Definieron el valor organizacional como «quién puede ejecutar de inmediato un experimento y explicar el failure». Esta cultura asegura que el code review cubra a la vez quality y compliance.

\subsection{Asignación de capital y portafolio de proyectos}

La inversión en GenAI no se apuesta enteramente en un solo agent, sino que se construye un portfolio de varias etapas:

\begin{itemize}
    \item Stage 1: GPT-style capability research continúa explorando scaling, architecture
    \item Stage 2: Platform / Agent ops brinda tooling + governance
    \item Stage 3: Conversión en producto convierte el agent en specific vertical solutions (como finance assistant)
\end{itemize}

\begin{importantbox}{Beneficios del portfolio por etapas}
Esta combinación permite que el equipo explore la innovación en el early-stage, garantice platform reliability en el mid-stage y asuma la commercialization en el late-stage. Los metrics / budget de cada etapa deben mantenerse claramente separados, para evitar que el departamento de capacidades y el departamento de product interfieran entre sí dentro del mismo sprint.
\end{importantbox}

\subsection{Resumen de la sección}

En el nivel organizacional es necesario establecer tres mecanismos clave —«densidad de talento», «squad multifuncionales» y «inversión por etapas»— para mantener al equipo bajo control mientras la tecnología evoluciona con rapidez.

\section{Repaso de las diapositivas clave}

\subsection{Portada y posicionamiento del tema}

\begin{figure}[H]
\centering
\includegraphics[width=0.8\textwidth]{cover.jpg}
\caption{Portada del curso: revisión de las tendencias de GenAI desde la perspectiva empresarial}
\end{figure}

La portada señala: esta clase no es una investigación de algoritmos, sino que aborda el tema desde cuatro niveles empresariales: GTM, plataforma, herramientas y gobernanza. Esto también determina que, en todo el conjunto de diapositivas, cada tema incluya cuatro dimensiones: casos, datos, cadena de herramientas y políticas.

\subsection{Mapa de tendencias y timeline}

En varias partes de las diapositivas aparece el mapa de tendencias: en el lado izquierdo se enumeran las capability técnicas (como multi-modal, agent), y en el lado derecho las necesidades empresariales (como compliance, observabilidad). Burak insiste repetidamente en colgar esta imagen en la pared, para que el equipo la consulte al elaborar el roadmap.

\begin{knowledgebox}{Cómo usar el mapa de tendencias}
\begin{enumerate}
    \item Primero ubica tu capability: ¿cuentas con un search + tool loop?
    \item Después revisa el lado derecho de las necesidades: ¿qué regulator, qué SLA aún no se han cumplido?
    \item Por último, dibuja el gap: el gap en sí mismo es la base de la prioritization
\end{enumerate}
\end{knowledgebox}

\subsection{Demostración de tiempo real y observabilidad}

Otra diapositiva clave muestra el dashboard de observabilidad de Vertex Agent: en el lado izquierdo está latency/time series, y en el lado derecho está tool-call matrix. Burak señala que el dashboard no es un nice-to-have, sino una de las condiciones del product launch.

\subsection{Resumen de la sección}

Las diapositivas no solo ofrecen un marco macro, sino que también pueden reutilizarse directamente en el workshop del equipo (mapa de tendencias + plantilla de dashboard).

\section{Análisis profundo del Agent Workflow}

\subsection{Descubrimiento y comprensión}

Cada agent empieza con intent detection: mapea la entrada en lenguaje natural a un structured goal. El equipo predefine slot templates para los intents comunes. La práctica que presenta Burak es usar primero un small model para hacer classification, y después un large model para hacer planning.

\subsection{Planeación y selección de herramientas}

El planner puede ser tanto un chain-of-thought prompt como una learned policy. Burak mostró un hybrid approach: primero se usa la policy para generar un plan tree, y después se usa el LLM para desarrollarlo paso a paso. Lo importante es el tool selection module: tiene que considerar el current status de la herramienta (si está sana), el estimated latency y el user context.

\subsection{Ejecución y revisión}

El Executor es, en realidad, una tool orchestration layer. Además de llamar a las API, tiene que manejar retries, fallbacks y confirm-with-user. Burak enfatiza que, para las operaciones financieras, siempre hay que agregar un explicit confirmation step, y registrar la decisión en el audit log.

\subsection{Cierre e iteración}

La última etapa del Workflow es la síntesis y el logging: se genera un summary para el usuario, se actualiza la memory, se registran todas las interacciones en el observability pipeline, y se dispara el training pipeline según el failure tag.

\subsection{Resumen de la sección}

Descomponer el agent workflow en cuatro pasos (descubrimiento, planeación, ejecución y cierre) ayuda a establecer métricas, monitoreo y revisión en cada etapa.
\section{Caso de gobernanza: playbook regulatorio}

\subsection{Playbook de regulación financiera}

Burak hace especial énfasis en los requisitos regulatorios del escenario financiero: cada recommendation debe ir acompañada de una fact sheet que cite explícitamente qué regulation aplica. El Agent pipeline genera automáticamente un `regulatory footprint` para la compliance review.

\subsection{Guardrail adicionales para el sector médico y el gobierno}

Las instituciones médicas exigen que el sistema ofrezca breach detection: antes de generar una recomendación de tratamiento, el agent primero verifica el patient consent y aplica la whitelist tool. Este tipo de guardrail suele combinar policy-as-code con real-time monitoring, formando un ciclo cerrado de «statement + enforcement».

\subsection{Resumen de la sección}

El playbook de gobernanza se compone de policy de distintos vertical; todas las policy se implementan mediante un gatekeeper convertido en code, lo que garantiza que las capacidades del agent no rebasen el alcance de las reglas.
\section{Apéndice: plantillas de Prompt y de Tool}

\subsection{Arquitectura del Prompt Template}

El prompt template que recomienda Burak consta de tres partes:

\begin{enumerate}
    \item Context: incluye los antecedentes de la tarea, las restricciones de policy y la intención del usuario
    \item Instruction: especifica los requisitos de salida del paso actual (por ejemplo, «enumera las top-3 soluciones»)
    \item Tool Usage: enumera los métodos de invocación de las herramientas disponibles, junto con un success indicator
\end{enumerate}

\begin{knowledgebox}{Lista de verificación del Prompt Template}
Cada vez que se actualiza el prompt, hay que verificar:
1. Si el Context se actualizó (policy / data)
2. Si la Instruction sigue siendo clara
3. Si las Tools \& examples son reproducibles
\end{knowledgebox}

\subsection{Ejemplo de Tool Manifest}

Un tool manifest típico incluye los siguientes campos:

\begin{itemize}
    \item tool name
    \item description
    \item input schema (tipos de parámetro, valores posibles)
    \item latency range
    \item safety constraints (qué user group no puede usarla)
\end{itemize}

Incluso si la tool es solo una API interna, se recomienda redactarla como manifest; el equipo de logística puede usarla como manual de operación.

\subsection{Sistema de etiquetas Failure Tag}

Burak recomienda unificar la failure tag taxonomy:

\begin{itemize}
    \item hallucination
    \item tool failure
    \item context expiration
    \item safety violation
    \item latency breach
\end{itemize}

Cada tag está asociado a un flujo de procesamiento: por ejemplo, `tool failure` activa fallback or escalate, y `hallucination` activa LLM retrace path.

\subsection{Resumen de la sección}

Las plantillas de prompt y de tool son la base de agent reproducibility; una vez estructurados los modos de fallo, el producto y la ingeniería pueden corregirlos con iteraciones más pequeñas.

\section{P\&A y temas de debate}

\subsection{Resumen de preguntas y respuestas del público}

En la segunda mitad de la entrevista, Burak respondió preguntas que le interesaban al público:

\begin{itemize}
    \item \textbf{P: ¿Cómo lograr que el agent se comporte de manera consistente en contextos multilingües?}\\
    R: Construir bien el pipeline de localización, internacionalizar tanto el prompt como el tool manifest, y usar una pila de embedding multilingüe.
    \item \textbf{P: ¿Cómo auditar la hallucination del agent?}\\
    R: Combinar search+LLM: al generar la answer, producir también la source cite correspondiente, y enviar tanto la cite como el user prompt al evaluation harness.
\end{itemize}

\subsection{Temas de debate y reflexión}

Burak planteó dos proposiciones que vale la pena debatir:

\begin{enumerate}
    \item ¿El guardrail de seguridad del agent debe ir primero o iterarse de manera simultánea?
    \item ¿Las capacidades Multi-modal deben desplegarse primero en task de bajo riesgo?
\end{enumerate}

Su postura es la siguiente: el guardrail necesita iterarse de manera simultánea (no se puede esperar a que la capacidad mejore para agregarlo); lo Multi-modal debe validarse primero en escenarios poco sensibles a la high-bandwidth observation, y luego extenderse gradualmente a escenarios de alta sensibilidad.

\subsection{Resumen de la sección}

La parte de P\&A reveló el pain point que más le interesaba al público, y también puso de relieve que las empresas deben mantener un debate constante entre la capacidad y la governance.

\section{Lista de verificación operativa}

\subsection{Deployment Readiness Checklist}

En la práctica de Burak, cada vez que un agent entra a producción se revisa la siguiente checklist:

\begin{itemize}
    \item El Metrics pipeline muestrea de forma continua Task Success, Tool Error y Latency
    \item El Observability dashboard ya está integrado con el monitoreo de DataDog / Vertex
    \item Las intercept tests de las Policy guardrails se aprueban
    \item El Feedback loop pipeline puede etiquetar automáticamente los failures y disparar el retrain
\end{itemize}

\subsection{Post-launch Monitoring}

Tras el lanzamiento, con el sprint como ciclo, cada semana se revisa:

\begin{itemize}
    \item ¿El alert channel tiene failures repetidos? Es necesario adjust thresholds
    \item si el new user feedback revela policy violation
    \item si la tool reliability (\#timeout, \#fallback) es estable
\end{itemize}

Solo el monitoreo continuo puede garantizar que el agent no tenga drift.

\subsection{Resumen de la sección}

La lista de verificación operativa fusiona la tecnología, la gobernanza y el monitoreo en un repeatable process, y es la clave para llevar las capacidades del laboratorio a producción.

\section{Perspectivas futuras y recomendaciones de acción}

\subsection{La ruta de los próximos 12 meses}

Burak predice que el próximo año se concentrará en dos direcciones: 1) unificar los agent en una single platform (dejar de usar un stack fragmentado); 2) mejorar la governance capability, especialmente la capacidad de auditoría de los multi-modal output.

\subsection{Recomendaciones estratégicas}

\begin{itemize}
    \item Convertir el agent capability roadmap en un multi-layer diagram (capability + policy + Ops)
    \item En early stage invest in tool manifest \& prompt atlas, en lugar de ir directamente al prompt final de GPT-4
    \item Adoptar `policy-as-code` tooling para que el compliance team participe en el sprint planning
\end{itemize}

\subsection{Resumen de la sección}

El próximo año se elevará la attention de capability a governance y Operations, y los equipos que se hayan preparado con anticipación ganarán en velocidad de implementación.

\section{Referencia rápida de términos}

\subsection{Tabla de términos clave}

\begin{tabular}{p{0.25\textwidth} p{0.7\textwidth}}
\toprule
Término & Explicación \\
\midrule
Agent Builder & Herramienta visual de Vertex AI para definir el planner, las tools y el feedback loop \\
Policy-as-code & Define las reglas de política mediante código ejecutable, combinándolo con OPA y otras herramientas para el gatekeeping \\
Tool manifest & Describe la entrada, la salida, la latency y el failure mode de cada herramienta \\
Feedback loop & Envía el user feedback y los failure tags al labeling \& retrain pipeline \\
Observability pipeline & Agrega latency, tool usage y alerts en el dashboard \\
\bottomrule
\end{tabular}

\begin{importantbox}{La función del glosario}
Internamente: ayuda a que los ingenieros en onboarding entiendan la plataforma rápidamente; de cara al exterior: le da al equipo de compliance definiciones unificadas. Se recomienda colocar el glosario en el README y darle mantenimiento periódico.
\end{importantbox}

\subsection{Aplicación y recomendaciones de mantenimiento}

\begin{itemize}
    \item Llevar el glossary bajo control de versiones y sincronizarlo con el README, para asegurar que los términos se actualicen en cuanto cambie la policy
    \item Durante el onboarding de los nuevos integrantes, pedirles que repitan el glossary para asegurar la comprensión y la memorización
    \item Convertir el glossary en un `cheatsheet` que se use en las compliance meeting y en el board review
\end{itemize}

\section{Material adicional e inspiración}

\subsection{Idea de combinación}

Al combinar el contenido de esta clase en un workshop, puede dividirse en tres bloques: 1) Trends Review, que empieza con un timeline y una trends table; 2) Platform Deep Dive, que aborda Vertex y Observabilidad; 3) Governance Sprint, que usa un prompt template y policy-as-code como hands-on exercise. Cada bloque puede tener una slide, para que el equipo repase todo el contenido en 3 horas.

\subsection{Citas e inspiración}

Burak mencionó varias veces que «la capacidad de observación importa más que un modelo nuevo»: mantener instrumentation, dashboarding y alerting te permite detectar a tiempo un regress cuando se actualiza el modelo, en lugar de esperar a que los usuarios se quejen. Esta idea también se refleja en la risk column de la Table 1: la observabilidad reduce directamente el compliance risk.

\subsection{Resumen de la sección}

El material adicional ayuda al equipo a convertir rápidamente el contenido de la talk en un workshop y un playbook, para que la narrative no se olvide en cuanto termine de escucharse, sino que realmente se implemente.

\section{Síntesis y ampliación}

\subsection{Tabla de síntesis global del curso}

\begin{table}[H]
\centering
\begin{tabular}{p{0.2\textwidth} p{0.23\textwidth} p{0.24\textwidth} p{0.23\textwidth}}
\toprule
Tema & Conclusión del curso & Acción clave & Riesgo / punto de atención \\
\midrule
Tendencias de GenAI empresarial & La democratización de la IA y su conversión en plataforma están ampliando la base de usuarios & Construir plataformas multimodales, preparar tooling elástico & Depender de un solo modelo o API genera fácilmente lock-in \\
Arquitectura del agente & search + LLM y la cadena de herramientas de Vertex AI son la base de la implementación & Refinar módulos como memory, retriever y planner & El costo de las fallas es alto y los errores en las llamadas a herramientas son frecuentes \\
Observabilidad y gobernanza & Las métricas, las alertas y el audit trail son un prerrequisito para el despliegue & Establecer una policy layer y un workflow human-in-loop & Sin un audit trail es difícil determinar responsabilidades \\
Operación del despliegue & El equilibrio entre latency y cost requiere una programación fina & multi-model routing, auto scaling y feedback loop & Sobredimensionar desperdicia el presupuesto y subdimensionar afecta la experiencia \\
Talento e innovación & Se necesita creatividad y una perspectiva interdisciplinaria & Invertir en la iteración de prompts, proyectos interdisciplinarios y valor social & Hacer solo investigación técnica hace perder de vista el impacto en el mercado y en la sociedad \\
\bottomrule
\end{tabular}
\caption{Conclusiones clave y acciones correspondientes de la Lecture 09}
\end{table}

\subsection{Acciones adicionales y reflexión}

\begin{itemize}
    \item Para los equipos de producto empresarial: convertir la combinación de search y LLM en un pipeline cerrado, e incorporar en él el monitoreo de calidad y el feedback loop
    \item Para los equipos de despliegue: convertir el uso de herramientas, la latency y el cost en un panel permanente, y hacer que las actualizaciones del modelo se ejecuten primero en un sandbox
    \item Para la gobernanza de seguridad: integrar la Policy Layer con el audit trail, para asegurar que, cuando el agente falle, se pueda rastrear rápidamente el origen
    \item Para investigadores y estudiantes: incorporar la creatividad y una perspectiva interdisciplinaria en los proyectos de IA, prestando atención al valor social y no solo a lo técnico
\end{itemize}

\subsection{Lecturas adicionales}

\begin{itemize}
    \item Google, ``Gemini: A Family of Highly Capable Multimodal Models,'' 2023.
    \item Google Cloud, ``Vertex AI Agents Documentation,'' 2024.
    \item Anthropic, ``Responsible Scaling Policy,'' 2023.
    \item Kamradt, ``Needle in a Haystack - Pressure Testing LLMs,'' 2023.
\end{itemize}

\end{document}
